\subsection{Context-Dependent Relevance: Calibrating Annotation Sources}

\begin{figure}[t!]
    \centering
    {{fig:model_selection}}
    \caption{Model selection across geneset libraries. (a) Performance of PIGEAN models defined by different geneset library combinations on common disease validation traits, measured by relative success and mechanistic alignment. Marker size indicates model complexity. (b) Same analysis on rare disease validation traits. (c) Proportion of significant genesets from each library across trait-model pairs.}
    \label{fig:model-selection}
\end{figure}

The concept of disease relevance is inherently context-dependent: a gene relevant in the context of mouse model studies may differ from one relevant to human therapeutic intervention. Our framework operationalizes this by allowing different geneset libraries to define different relevance contexts. A model trained on mouse knockout phenotypes captures relevance as defined by mammalian loss-of-function studies; a model incorporating all available annotations captures a broader notion of relevance anchored by the full scope of current biological knowledge; and a simplified model selected for performance across diverse applications provides a practical balance of accuracy and efficiency.

To evaluate how annotation sources contribute to indirect support, we constructed multiple PIGEAN models, each defined by a specific combination of geneset libraries, and assessed their performance on held-out validation sets (Figure~\ref{fig:model-selection}a-b). Libraries included knockout mouse phenotypes ({{NUM_MPO_TERMS:,}} Mammalian Phenotype Ontology terms), curated pathways ({{NUM_MSIGDB_GENESETS:,}} MSigDB gene sets), literature co-occurrence ({{NUM_PFOCR_SETS:,}} Pathway Figure OCR sets, {{NUM_MESH_TERMS:,}} MeSH terms), protein interactions ({{NUM_STRING_NEIGHBORS:,}} STRING neighbors), and single-cell expression patterns ({{NUM_SINGLE_CELL_CLUSTERS:,}} cell type clusters).

% TODO: Add specific model comparison results
% Which combinations performed best on common vs rare?
% How much does adding more libraries help vs. hurt?

Across model selection experiments, a parsimonious model combining mouse phenotypes and MSigDB gene sets (Mouse+MSigDB) emerged as the top performer for general-purpose applications (Figure~\ref{fig:model-selection}a-b). This model achieved high relative success for drug target prioritization while maintaining mechanistic interpretability, and we adopted it for phenome-wide analysis. The spike-and-slab prior inherent to PIGEAN regularizes away uninformative annotations, enabling robust performance even when including heterogeneous or partially redundant libraries.

The contribution of each library varied systematically across trait categories (Figure~\ref{fig:model-selection}c). Mouse phenotype annotations contributed disproportionately to metabolic and developmental traits, consistent with the extensive phenotyping of these systems in model organisms. MSigDB pathway annotations showed broader utility across trait categories. Literature-derived annotations (PFOCR, MeSH) contributed more to well-studied diseases but showed signatures of ascertainment bias toward historically researched genes. These patterns highlight that the ``best'' annotation source depends on both the trait of interest and the intended application.

% ANALYSIS IDEA: Show that model selection is predictive - the best model for a trait
% category on the validation set also performs best on held-out traits in that category.

% TODO: Add quantitative results for Mouse+MSigDB model selection
% - How many traits converged?
% - What was the average number of significant genesets per trait?
% - How did complexity (number of genesets) relate to performance?
